\documentclass[11pt]{article}
\usepackage[margin=1in]{geometry}
\usepackage{graphicx}
\usepackage{amsmath, amssymb}
\usepackage{booktabs}
\usepackage{hyperref}
\usepackage{url}
\usepackage{natbib}
\bibliographystyle{plainnat}

\title{A multiscale computational framework linking biochemical pattern
formation to tissue-level spine morphogenesis in molluscan shells}
\author{}
\date{}

\begin{document}
\maketitle

\begin{abstract}
Connecting molecular-scale processes to organism-level morphology is a
central challenge in computational biology, particularly in non-model
organisms for which detailed gene regulatory knowledge is lacking. Existing
frameworks for molluscan shells describe either biochemical pigmentation
patterning or tissue-scale mechanics, but rarely bridge the two. We present
a modular multiscale framework that couples a two-species Activator-Substrate
reaction-diffusion model at the molecular scale to a morphoelastic rod
description of mantle growth at the tissue scale. A Quality-Diversity
algorithm (MAP-Elites, implemented in the \texttt{qdpy} Python library)
builds a pre-computed archive indexed by the parameters $\{g_1, g_2\}$ of a
Gaussian growth profile, allowing rapid inference of molecular parameters
from tissue-scale fits. We apply the framework to the gastropod
\textit{Turbo sazae} using specimens collected from three Japanese
localities: $n=12$ from Hayama, $n=16$ from Jyogashima, and $n=9$ from
Tateyama. Two-sided Wilcoxon rank-sum tests show that the distribution of
$g_1$ values for Jyogashima shells is on average $0.25$ higher than for
Hayama ($p = 1.202 \times 10^{-5}$) and $0.4$ higher than for Tateyama
($p = 5.7 \times 10^{-6}$), while the two sheltered-shore populations do not
differ significantly. The corresponding biochemical parameters shift in a
direction consistent with increased activator production and reduced
degradation. These results quantify a previously unreported phenotypic
difference between populations and, more broadly, illustrate how a
multiscale framework can trace morphological variation back to candidate
molecular mechanisms in a non-model organism.
\end{abstract}

\section{Introduction}

One of the enduring goals in computational biology is to link quantitative
descriptions across disparate spatial and temporal scales, from molecular
interactions to organism-level form. Advances in each scale in isolation
have been substantial, but the construction of unified frameworks that
explicitly bridge the molecular and tissue scales remains a fundamental
challenge \citep{goriely2017mathematics}. In the context of growth and
morphogenesis, tissue-level mechanical models of continuum media have been
used to great effect \citep{moulton2013morphoelastic,moulton2020morphoelastic},
while biochemical models drawing on reaction-diffusion theory have proven
remarkably flexible at reproducing diverse spatial patterns
\citep{turing1952chemical,gierer1972theory,meinhardt1987model,pueyo2018reactiondiffusion}.
The coupling of these scales, however, has received comparatively little
attention even in systems where the natural connection is especially clear.

Molluscan shells provide a compelling study system for cross-scale modelling
of pattern formation. The shell preserves a spatiotemporal record of its
own secretion by the mantle, and the iterative geometry of this secretion
has attracted quantitative attention since the foundational work of Raup
\citep{raup1961geometric,raup1966geometric}. Geometric models have since
been extended and reviewed in subsequent literature
\citep{gerber2017geometry,urdy2011recapitulation,meinhardt1995algorithmic},
but the parameters of these descriptions are typically disconnected from
the underlying biological processes of shell secretion. A parallel strand
of research has explained shell \emph{ornamentation}, including spines and
ribs, through the mechanics of a soft mantle attached to a rigid shell
aperture \citep{chirat2013mechanical,moulton2015morpho,erlich2016morphomechanics}.
These mechanics-based accounts typically require a spatial pre-pattern in
growth or stiffness as input, yet they remain silent on where that
pre-pattern comes from. Biochemical descriptions of shell pigmentation,
meanwhile, have been developed largely independently
\citep{fowler1992modeling,meinhardt1995algorithmic,boettiger2009neural,gierer1972theory}.
Despite the natural suggestion that structural pattern formation may be
governed by an underlying chemical pre-pattern, no prior framework has, to
our knowledge, explicitly coupled biochemical pigmentation-style models to
tissue-level mechanical models of ornamentation.

In this work, we address that gap by constructing a multiscale computational
framework that links a two-species Activator-Substrate reaction-diffusion
model at the molecular scale to a morphoelastic rod description of mantle
growth at the tissue scale. To keep the framework tractable, we treat the
molecular model as a generator of spatially structured growth profiles which
are then passed to the mechanical model as input. Because the tissue-scale
model is expensive, we pre-integrate it over a grid of growth parameters
and use a Quality-Diversity algorithm, MAP-Elites
\citep{mouret2015illuminating,cully2015robots}, implemented with
\texttt{qdpy} \citep{cazenille2018qdpy}, to fill an archive that maps the
molecular parameter space to the tissue-scale morphospace. The archive then
functions as a look-up table, bypassing the mechanical bottleneck during
inference. Figure~\ref{fig:framework} summarises the framework.

\begin{figure}[htbp]
\centering
\includegraphics[width=0.95\textwidth]{figures/fig_framework.pdf}
\caption{Schematic of the multiscale framework. A molecular Activator-Substrate
reaction-diffusion model with parameters $\{c_1,c_2,c_3,c_4,d_1,d_2\}$
generates a spatial concentration profile that is mapped to the parameters
$\{g_1, g_2\}$ of a Gaussian growth profile at the tissue scale. The
morphoelastic rod model then evolves a spine shape which is compared against
real shell images. A MAP-Elites archive connects the tissue-scale
$\{g_1, g_2\}$ space back to molecular parameter sets.}
\label{fig:framework}
\end{figure}

As a testbed, we apply the framework to the gastropod \textit{Turbo sazae}
\citep{fukuda2017nomenclature}, a species native to coastal Japan that
exhibits remarkable plasticity in spine form: populations from exposed
shores often display prominent horns, while those from sheltered bays can
develop only short horns, or none at all
\citep{tanaka2018phenotypic,hamm2008environmental}. We analyse shells
collected from three Japanese localities with distinct environmental
characteristics and ask whether any morphological difference in spine shape
can be detected and, if so, whether the multiscale framework can trace that
difference back to parametric shifts at the biochemical level.

Our contributions are fourfold. \textbf{(i)} We build a modular multiscale
framework that explicitly maps biochemical parameter sets to tissue-level
spine shapes through an intermediate Gaussian growth profile.
\textbf{(ii)} We show how pre-integration of the mechanical model, coupled
with MAP-Elites exploration, turns the framework into a practical inference
procedure. \textbf{(iii)} We apply the framework to three populations of
\textit{T. sazae} and uncover a statistically significant difference in
spine form, with Jyogashima shells exhibiting $g_1$ values on average
$0.25$ higher than Hayama ($p = 1.202 \times 10^{-5}$) and $0.4$ higher
than Tateyama ($p = 5.7 \times 10^{-6}$) by two-sided Wilcoxon rank-sum
tests. \textbf{(iv)} We trace that phenotypic shift through the framework
to higher values of the biochemical parameter $c_2$ and lower values of
$c_3$ in the Jyogashima population, providing a concrete molecular
candidate for environmentally driven phenotypic variation. The remainder of
the paper details the modelling components (Section~\ref{sec:methods}),
presents the morphospace and population-level findings
(Section~\ref{sec:results}), and discusses implications and limitations
(Section~\ref{sec:discussion}).

\section{Background and Related Work}

\paragraph{Theoretical morphology of shells.} The mathematical modelling of
shell form traces back to the seminal morphospace parameterization of the
coiled shell by Raup \citep{raup1961geometric,raup1966geometric}, and has
since been extended by a wide literature on geometric and generative models
\citep{gerber2017geometry,urdy2011recapitulation,meinhardt1995algorithmic}.
These approaches describe form through a small set of geometric parameters
but do not attempt to ground those parameters in cellular or subcellular
processes.

\paragraph{Mechanical models of shell ornamentation.} A parallel line of
work treats the mantle-shell interface as a growing elastic structure
subject to mechanical buckling when a length mismatch develops. Chirat et
al.\ \citep{chirat2013mechanical} demonstrated that spine morphologies on
molluscan shells can be explained by the elastic buckling of a mantle that
has grown longer than the current shell aperture. This idea has since been
extended to ammonites and rib patterns
\citep{moulton2015morpho,erlich2016morphomechanics}. A recurring feature is
that generating realistic shapes requires a spatial pre-pattern in growth
or mechanical properties as an input; the origin of that pre-pattern is not
resolved within mechanics alone.

\paragraph{Biochemical pattern formation.} Reaction-diffusion dynamics
underpin most biochemical models of shell pigmentation. Following Turing's
chemical theory of morphogenesis \citep{turing1952chemical}, Gierer and
Meinhardt introduced the activator-inhibitor architecture
\citep{gierer1972theory}, and Meinhardt and Klingler
\citep{meinhardt1987model} applied reaction-diffusion modelling to
molluscan shells, forming the basis of subsequent work on shell pattern
generation \citep{fowler1992modeling,meinhardt1995algorithmic}. Boettiger
et al.\ \citep{boettiger2009neural} introduced a neural-inspired
formulation. The two-species Activator-Substrate variant used in our
framework is inherited from this lineage.

\paragraph{Quality-Diversity optimisation.} MAP-Elites
\citep{mouret2015illuminating} is a black-box optimisation algorithm that
maintains a structured archive of diverse, high-performing solutions. The
algorithm has found broad application in evolutionary robotics
\citep{cully2015robots} and is implemented in the \texttt{qdpy} Python
library \citep{cazenille2018qdpy}. Here we use it to fill a
two-dimensional archive indexed by the growth-profile parameters
$\{g_1, g_2\}$ with molecular parameter sets whose concentration profile
best approximates a Gaussian of those parameters.

\paragraph{Phenotypic plasticity in molluscs.} \textit{Turbo sazae} is a
marine gastropod previously conflated taxonomically with \textit{Turbo
cornutus}; the two were separated in the nomenclatural revision of Fukuda
\citep{fukuda2017nomenclature}. The species is known for plasticity in
horn form. More broadly, shell plasticity in gastropods has been reviewed
in the context of ecology and evolution
\citep{tanaka2018phenotypic,hamm2008environmental}, but quantitative
morphometric analyses have typically not been connected to underlying
biochemical processes.

\section{Methods}
\label{sec:methods}

\subsection{Molecular-level Activator-Substrate model}

At the molecular scale, we adopt the Activator-Substrate reaction-diffusion
system from the lineage of \citet{fowler1992modeling}, in which an
activator concentration $a(x,t)$ and a substrate concentration $s(x,t)$
interact autocatalytically along a one-dimensional spatial coordinate.
The non-dimensionalised dynamics are governed by six parameters,
$S := \{c_1, c_2, c_3, c_4, d_1, d_2\}$. The parameter $c_1$ controls
saturation of the autocatalytic production, $c_2$ governs direct production
of activator from substrate, $c_3$ is the activator degradation rate, $c_4$
is the substrate decay rate, and $d_1$ and $d_2$ are the diffusion
coefficients of the two species. Starting from a slightly perturbed
spatially uniform initial condition, the system can settle into stable
spatial profiles of the activator species. These profiles are the output
that we pass to the tissue-level model.

The diffusion coefficients are fixed throughout this study, because we do
not expect variation in diffusion on ecological timescales. The remaining
four kinetic parameters, $\{c_1, c_2, c_3, c_4\}$, are sampled by the
Quality-Diversity algorithm. We follow the interpretation in
\citet{meinhardt1987model} of the activator-substrate framework as a
phenomenological stand-in for an unknown gene regulatory network. This is
a deliberate modelling choice rather than a mechanistic commitment: the
molecular model is modular and can in principle be replaced by more
realistic biology without altering the overall architecture of the
framework.

\subsection{Tissue-level morphoelastic rod model}

At the tissue scale we follow the line of work on shell mechanics initiated
by \citet{chirat2013mechanical} and formalised in the morphoelastic rod
formalism of \citet{moulton2013morphoelastic}. The mantle edge is treated
as a planar elastic rod parameterised by an initial arc length coordinate
$S_0$ and time $t$, with position $\mathbf{r}(S_0, t)$. The rod evolves
quasistatically under the vector equilibrium conditions for the resultant
force $\mathbf{n}$, resultant moment $\mathbf{m}$, and external force
$\mathbf{f}$ imposed by the shell edge. For a planar rod, the moment and
curvature take the form $\mathbf{m} = m \mathbf{e}_z$ and
$\mathbf{u} = u \mathbf{e}_z$, and are related by a linear constitutive
law, $m = E_b u$, with $E_b$ the bending stiffness.

The shell edge, parameterised by $\mathbf{p}(S_0, t)$, interacts with the
mantle through the adhesive force $\mathbf{f} = k(\mathbf{p} - \mathbf{r})$,
where $k$ sets the strength of the attachment. Calcification updates the
shell edge via an evolution equation with rate constant $\eta$, and mantle
growth is described within the morphoelastic framework
\citep{moulton2013morphoelastic,goriely2017mathematics} by a growth stretch
$\gamma = \partial S / \partial S_0$ that evolves according to a growth rule
$g$. Critically, $g$ may depend on position, introducing the spatial
heterogeneity that is needed for realistic spine formation
\citep{chirat2013mechanical}.

To keep inference tractable, we restrict the incremental growth profile to
a Gaussian form parameterised by two scalars,
$g_1$ (amplitude) and $g_2$ (width),
\begin{equation}
g(s) \propto g_1 \, \exp\!\left(-\frac{(s - s_0)^2}{2 g_2^2}\right),
\end{equation}
where $s$ is the current arc-length coordinate and $s_0$ is the centre of
the growing region. Larger $g_1$ corresponds to a higher incremental growth
rate, and larger $g_2$ to a wider growing region. Bending stiffness is
linked directly to growth, producing realistic spine shapes more naturally
than in prior decoupled models. Domain width, calcification rate, and
adhesion stiffness are set in non-dimensional form by combining physical
estimates with the requirement that realistic buckling modes emerge.

\subsection{Quality-Diversity exploration}

The key computational challenge is that solving the tissue-level model for
a single parameter set takes on the order of minutes, which makes any
direct search of the six-dimensional molecular parameter space infeasible.
We address this in two steps. First, we pre-integrate the tissue-level
model over a two-dimensional grid of $\{g_1, g_2\}$ values to produce a
finite zoology of time-evolving spine shapes. Second, we use MAP-Elites
\citep{mouret2015illuminating,cully2015robots}, implemented in \texttt{qdpy}
\citep{cazenille2018qdpy}, to fill an archive indexed by $\{g_1, g_2\}$
with molecular parameter sets whose activator concentration profile best
approximates a Gaussian of those parameters. The quality criterion is the
goodness-of-fit of the molecular profile to a Gaussian, and the diversity
criterion is the cell location in the $\{g_1, g_2\}$ grid. Generation of
new candidate parameter sets proceeds by random sampling or mutation of
existing archive entries, with each parameter sampled either linearly or
logarithmically according to its range. The exploration continues until a
specified evaluation budget is exhausted, and multiple independent runs
are combined to improve archive coverage.

\subsection{Shell collection, imaging, and parameter extraction}

We collected \textit{T. sazae} shells from three Japanese localities: Najima
in Hayama and Jyogashima in Miura, both in Kanagawa Prefecture, and
Okinoshima in Tateyama, Chiba Prefecture. Hayama and Tateyama are
relatively sheltered shores with ongoing coastal desertification, while
Jyogashima is a more exposed locality that has been recovering from
desertification and offers greater food availability. Specimens were
obtained from local fishermen and dissected to extract the shells on the
day of capture. Shells were photographed with a digital camera (Canon EOS
Kiss x7) equipped with a macro lens (Canon EF100mm f/2.8L Micro IS USM),
and spines were imaged in magnified view using a Nikon SMZ10 dissection
microscope mounted with the same digital camera. The final dataset
comprises $n=12$ shells from Hayama, $n=16$ from Jyogashima, and $n=9$
from Tateyama. Figure~\ref{fig:samples} summarises these sample sizes.

For each shell, we extracted best-fit tissue-scale parameters by overlaying
model-generated spine profiles onto the magnified spine photograph inside
a Mathematica GUI. Model profiles were first scaled, rotated, and
translated to match the size and orientation of the real spine. The
parameters $\{g_1, g_2\}$ were then adjusted while stepping through the
time sequence of simulated profiles, with early and middle growth lines in
the real spine compared against the corresponding simulated stages. Small
spines that had not accumulated enough structural detail were excluded, as
almost any parameter set could fit a shallow bump equally well. All shells
were fit in sequence before any population-level analysis to avoid
confirmation bias.

\begin{figure}[htbp]
\centering
\includegraphics[width=0.6\textwidth]{figures/fig_sample_sizes.pdf}
\caption{Sample sizes of \textit{Turbo sazae} shell specimens collected at
each of the three Japanese localities: $n=12$ from Hayama, $n=16$ from
Jyogashima, and $n=9$ from Tateyama.}
\label{fig:samples}
\end{figure}

\subsection{Population-level statistical comparison}

We compared the distributions of $g_1$ and $g_2$ across the three
populations using two-sided Wilcoxon rank-sum tests. The test was chosen
to avoid distributional assumptions given the modest per-population sample
sizes and the discrete, grid-induced structure of the extracted $\{g_1,
g_2\}$ values. Biochemical parameters $c_1, \dots, c_4$ associated with
each population were then aggregated across MAP-Elites runs and summarised
by their population-wise mean.

\section{Results}
\label{sec:results}

\subsection{Growth parameters shape the spine morphospace}

Before turning to the shell data, we examined how the growth-profile
parameters $\{g_1, g_2\}$ shape the simulated spine morphology. The
context for this analysis is that spine shape depends on the incremental
growth profile via the morphoelastic dynamics of the mantle; understanding
how each parameter influences shape is a prerequisite for interpreting any
subsequent differences between populations. We therefore pre-integrated
the tissue model on a grid covering plausible values of $g_1$ and $g_2$
and inspected the resulting morphospace.

Two trends are apparent. Increasing $g_1$, which controls the rate of
incremental growth, produces a more curved spine that reaches the point of
self-contact at a lower height. Intuitively, a larger excess of length per
time step drives a more pronounced mantle deformation. Conversely, smaller
$g_1$ yields a straighter, taller spine before self-contact occurs,
because the mantle deforms less aggressively per step and stays closer to
the evolving shell edge. Increasing $g_2$, which controls the width of the
growing region, produces a wider spine overall. A narrower growing region
additionally sharpens the spine tip, an effect amplified by the coupling
between bending stiffness and growth. Figure~\ref{fig:morphospace}
illustrates these trends schematically.

Taken together, these observations confirm that $\{g_1, g_2\}$ provides a
two-dimensional control over the most salient qualitative features of
spine shape, namely tip curvature and basal width. This justifies the
restriction to a Gaussian growth profile at the tissue scale: two degrees
of freedom are sufficient to produce a diverse morphospace of realistic
shapes while keeping the downstream inference problem low-dimensional.

\begin{figure}[htbp]
\centering
\includegraphics[width=0.65\textwidth]{figures/fig_morphospace.pdf}
\caption{Qualitative morphospace of simulated spine shapes as the Gaussian
growth parameters $g_1$ (rate) and $g_2$ (width) are varied. Larger $g_1$
produces more curved spines that reach self-contact earlier, while larger
$g_2$ produces wider spines. Schematic only; the axes illustrate the
structural trend described in the text.}
\label{fig:morphospace}
\end{figure}

\subsection{Biochemical parameters map onto the growth morphospace}

Having established the tissue-scale morphospace, we next asked which
regions of the molecular parameter space can generate Gaussian-like
concentration profiles corresponding to given $\{g_1, g_2\}$ pairs. The
MAP-Elites exploration produced an archive in which each cell of the
$\{g_1, g_2\}$ grid stores the biochemical parameter set whose activator
profile best matches a Gaussian of those parameters.

The aggregated archive shows a clear global organisation. The parameters
$c_1$, $c_2$, and $c_3$ each carve out distinct regions in the $\{g_1,
g_2\}$ plane. Low $c_1$ (weak autocatalytic saturation), combined with low
$c_2$ (weak direct production) and high $c_3$ (fast degradation), yields
low $g_2$ values; high $c_1$ is associated with high $g_2$ values; and
intermediate $g_2$ values arise when $c_2$ is high with low $c_1$ and
$c_3$, consistent with strong total activation from both autocatalysis and
direct production. The parameter $c_4$, in contrast, is poorly localised
and spans its full range across the archive. This asymmetry is a
straightforward consequence of the biochemical model: $c_1$, $c_2$, and
$c_3$ all enter the activator equation directly, whereas $c_4$ appears
only in the substrate equation and influences the activator indirectly.

The archive has well-defined borders, with sparse populations in some
peripheral regions and a sharp absence of entries in others. The sharpness
of the upper-right boundary is consistent with a standard linear stability
analysis of the reaction-diffusion system: high values of both $g_1$ and
$g_2$ simultaneously correspond to a region where the uniform state is
stable and no spatial pattern emerges. Ruggedness at the lower edge of the
main cluster is more likely an artefact of finite sampling and could be
refined with targeted exploration. These observations together confirm
that the molecular model can access most but not all of the tissue-scale
morphospace, and that the boundary is itself informative about the
dynamics of the molecular model.

\subsection{Population-level differences in spine form}

We next applied the framework to the three populations of
\textit{T. sazae}. For every shell we extracted the best-fit
$\{g_1, g_2\}$ pair via the overlay procedure described in the Methods,
then compared the resulting distributions across populations using
two-sided Wilcoxon rank-sum tests.

A clear population-level bias emerged in $g_1$. Spines from shells
collected in Hayama and Tateyama are better fit by smaller values of
$g_1$, while spines from Jyogashima shells favour larger $g_1$ values and
additionally show a wider distribution. Quantitatively, the distribution
of $g_1$ values for Jyogashima is on average $0.25$ higher than for Hayama
with a $p$-value of $1.202 \times 10^{-5}$, and $0.4$ higher than for
Tateyama with a $p$-value of $5.7 \times 10^{-6}$
(Figure~\ref{fig:g1shift}).

\begin{figure}[htbp]
\centering
\includegraphics[width=0.7\textwidth]{figures/fig_population_g1.pdf}
\caption{Population-level differences in the tissue-scale growth parameter
$g_1$. Two-sided Wilcoxon rank-sum tests show that $g_1$ values for the
Jyogashima population are on average $0.25$ higher than those from Hayama
($p = 1.202 \times 10^{-5}$) and $0.4$ higher than those from Tateyama
($p = 5.7 \times 10^{-6}$). Hayama and Tateyama do not differ significantly.}
\label{fig:g1shift}
\end{figure} The two sheltered-shore populations, Hayama
and Tateyama, do not differ significantly in $g_1$, and comparable
analysis on $g_2$ yields only very small differences between any pair of
populations. This implies that the single axis along which Jyogashima
shells differ from the other two is the one controlling spine curvature
and self-contact height.

Connecting this finding back to the molecular scale, the biochemical
parameters $c_1, \dots, c_4$ aggregated through the MAP-Elites archive
show a population-level shift consistent with the change in $g_1$. On
average, Jyogashima spines correspond to higher $c_2$ and lower $c_3$
than Hayama and Tateyama spines, while $c_1$ and $c_4$ differ only
weakly. This is mechanistically coherent: $c_2$ directly produces
activator from substrate and $c_3$ degrades it, so higher $c_2$ combined
with lower $c_3$ raises the steady-state activator concentration, which
in turn increases the tissue-scale growth rate $g_1$ and shifts the spine
into the more curved regime. No comparable trend appears in $c_4$, and
the stability analysis confirms that, as long as $c_4$ remains small
enough, its influence on the Gaussian profile is limited.

Interpreted together, these results quantify a previously unreported
morphological difference between a single \textit{T. sazae} population
and the other two we sampled, and they identify a concrete molecular
candidate for that difference. The statistical signal is strong
($p < 10^{-5}$ in both comparisons) despite the modest per-population
sample sizes, and it is consistent with the mechanistic direction
predicted by the model.

\begin{table}[htbp]
\centering
\caption{Sample sizes and statistical comparison of $g_1$ distributions
across populations. Comparisons against Jyogashima use the two-sided
Wilcoxon rank-sum test.}
\label{tab:populations}
\begin{tabular}{lccc}
\toprule
Population & $n$ shells & Mean $g_1$ shift vs.\ Jyogashima & $p$-value \\
\midrule
Hayama     & 12 & $-0.25$ & $1.202 \times 10^{-5}$ \\
Jyogashima & 16 & $0$ (reference) & --- \\
Tateyama   & 9  & $-0.40$ & $5.7 \times 10^{-6}$ \\
\bottomrule
\end{tabular}
\end{table}

\begin{table}[htbp]
\centering
\caption{Summary of model parameters at the molecular and tissue scales.}
\label{tab:parameters}
\begin{tabular}{lll}
\toprule
Scale & Parameter & Role \\
\midrule
Molecular & $c_1$ & Saturation of autocatalytic activator production \\
Molecular & $c_2$ & Direct activator production from substrate \\
Molecular & $c_3$ & Activator degradation rate \\
Molecular & $c_4$ & Substrate decay rate \\
Molecular & $d_1$ & Activator diffusion coefficient (fixed) \\
Molecular & $d_2$ & Substrate diffusion coefficient (fixed) \\
Tissue    & $g_1$ & Gaussian growth amplitude (rate) \\
Tissue    & $g_2$ & Gaussian growth width \\
\bottomrule
\end{tabular}
\end{table}

\section{Discussion}
\label{sec:discussion}

\subsection{Summary of findings}

The central result of this work is a multiscale computational framework
that links biochemical pattern formation to tissue-level mechanics in a
way that is both modular and empirically testable, applied to the
spine morphology of \textit{T. sazae}. Applied to three Japanese
populations, the framework uncovers a statistically significant
population-level shift in the tissue-scale growth parameter $g_1$: spines
from Jyogashima shells have $g_1$ values on average $0.25$ higher than
those from Hayama ($p = 1.202 \times 10^{-5}$) and $0.4$ higher than
those from Tateyama ($p = 5.7 \times 10^{-6}$). Tracing that shift
through the framework points to higher values of the biochemical
parameter $c_2$ and lower values of $c_3$ in the Jyogashima population.
The framework thereby provides not only a quantitative measurement of a
previously undescribed morphological difference, but also a parametric
hypothesis about its molecular origin.

\subsection{Connecting parameters to environment}

A natural next step is to link the inferred parametric differences to
environmental differences between the three localities. This is not
possible to do rigorously here, because our biochemical model is
deliberately phenomenological and is not calibrated against measured
molecular concentrations. Nevertheless, a consistent narrative emerges.
Jyogashima is an exposed shore that has been recovering from coastal
desertification and has greater seaweed availability, whereas Hayama and
Tateyama are sheltered sites that continue to suffer from desertification.
The inferred direction of the biochemical shift, namely higher $c_2$ and
lower $c_3$ in Jyogashima, translates to a higher effective growth rate
of the spine-producing tissue. It is plausible that greater food
availability translates through the metabolism of the animal into
elevated rates of molecular synthesis and reduced degradation. Wave
exposure is a second candidate environmental variable, though it is
confounded with food availability across our three sites. Disentangling
these factors is an empirical question that this work cannot settle but
that our framework makes concrete and quantitative.

\subsection{Limitations}

Several simplifying assumptions qualify the interpretations above. First,
the biochemical model is phenomenological; it is adapted from the
pigmentation lineage \citep{fowler1992modeling,meinhardt1987model} and
does not correspond to a measured gene regulatory network for spine
formation in this species or any close relative. Second, we model a
single isolated spine and do not account for the spatially extended,
temporally recurrent production of multiple spines across a full shell.
Third, we restrict the tissue-scale growth profile to a Gaussian form with
two degrees of freedom; more complex profiles might be needed to capture
the full diversity of spine shapes, particularly the multi-lobed
morphologies observed in some molluscan clades
\citep{chirat2013mechanical}. Fourth, the archive exploration is finite,
and the coverage of the lower edge of the morphospace could be improved
with biased or adaptive sampling. These limitations mean that the
biochemical parameter values we infer should be interpreted as directional
rather than absolute.

\subsection{Relation to prior work}

The closest antecedent at the tissue scale is the spine-formation model of
Chirat et al.\ \citep{chirat2013mechanical}, which used an inhomogeneous
bending stiffness combined with a uniform growth profile to generate
realistic spines. We retain the underlying morphoelastic machinery but
invert the architecture: growth is inhomogeneous, set by the molecular
activator profile, and bending stiffness is linked directly to growth.
This tighter coupling reflects an intuition that sharp spine tips
(high curvature) and soft mantle (low stiffness) should emerge from a
common upstream driver rather than from two independent spatial fields.
At the biochemical end, our model inherits the activator-substrate
vocabulary of pigmentation studies
\citep{fowler1992modeling,meinhardt1987model} and extends it from
colour patterning to structural patterning. MAP-Elites
\citep{mouret2015illuminating} has to our knowledge not previously been
used to search the parameter space of a reaction-diffusion model coupled
to a mechanical model.

\subsection{Future directions}

The framework is modular at each scale. The most immediate extension is
to replace the phenomenological Activator-Substrate module with a
biologically grounded network, if and when such a network becomes
available for spine-forming molluscs. Transcriptomic profiling of the
mantle, either bulk or at single-cell resolution, is a promising route
to generate candidate networks. Beyond the molecular module, the tissue
model can be extended to a whole-shell setting, where multiple spines
appear sequentially along the growing aperture, and to three dimensions.
At the organism level, applying the framework to additional populations
of \textit{T. sazae} across a wider environmental gradient would test the
generality of the parametric trend we observed in Jyogashima. Finally,
the framework itself is not specific to molluscan spines: any system
displaying morphological variation driven by a combination of
heterogeneous mechanics and an upstream biochemical signal is a
candidate, including leaf teeth, horns, and anatomical features shaped by
growth-coupled buckling. The role of this paper is to provide both a
concrete instantiation and a template for those extensions.

\bibliography{references}

\end{document}
